\begin{table}[!htbp]\centering
\caption{Large and small peers as sources of information}\label{tab:bigsmall}
\footnotesize
\begin{tabular*}{\linewidth}{@{\extracolsep{\fill}}lcccccc}
\toprule
 & \multicolumn{2}{c}{All firms} & \multicolumn{2}{c}{Large firms} & \multicolumn{2}{c}{Small firms} \\
\cmidrule(lr){2-3}\cmidrule(lr){4-5}\cmidrule(lr){6-7}
 & (1) & (2) & (3) & (4) & (5) & (6) \\
\midrule
Large-cap peers' return & 0.056 & 0.064 & 0.095** & 0.028 & 0.040 & 0.091 \\
 & (1.17) & (1.46) & (2.38) & (0.48) & (0.62) & (1.22) \\
Small-cap peers' return & 0.058*** & 0.090*** & 0.025 & 0.044 & 0.084*** & 0.133** \\
 & (3.10) & (3.21) & (1.54) & (1.28) & (3.00) & (2.58) \\
\addlinespace
Small $-$ large & 0.002 & 0.026 & -0.070 & 0.016 & 0.044 & 0.043 \\
 & (0.04) & (0.46) & (-1.56) & (0.24) & (0.62) & (0.52) \\
\addlinespace
Controls &  & Yes &  & Yes &  & Yes \\
Observations & 13,197 & 10,689 & 6,622 & 5,237 & 6,575 & 5,452 \\
\bottomrule
\end{tabular*}
\par\vspace{4pt}\begin{minipage}{\linewidth}\scriptsize \emph{Notes:} Each firm's customer-overlap peers are split at the month-$t$ median market capitalization of sample firms; each sub-portfolio is weighted by shared customers. Fama-MacBeth regressions of month-$t+1$ returns for all firms and for firms above and below the median size, without controls (columns (1), (3), and (5)) and with the controls of Table~\ref{tab:fm}, column (5) (columns (2), (4), and (6)). The sample is firms with at least one customer-overlap peer; when one sub-portfolio is empty, its return is set to the monthly cross-sectional mean. ``Small $-$ large'' is the time-series average of the monthly difference between the two slopes. Newey-West $t$-statistics (four lags) in parentheses. *, **, and *** denote significance at the 10\%, 5\%, and 1\% levels.\end{minipage}
\end{table}